%=====================================================
\section{Tier 1.1: Mathematical formulation}
\label{sec:tier1_1}

This section summarizes the governing equations corresponding to the input file
\texttt{tier1\_1\_mode\_I\_monotonic.i} and the analytical reference implemented in
\texttt{verify\_tier1\_1.py}.

%=====================================================
\subsection{Domain and discontinuity}

We consider two stacked unit-square subdomains,
\begin{equation}
\Omega_1 = [0,1]\times[0,1],
\qquad
\Omega_2 = [0,1]\times[1,2],
\end{equation}
with the interface
\begin{equation}
\Gamma_c = \{(x,y): 0 \le x \le 1,\ y = 1\}.
\end{equation}

The mesh is split along $\Gamma_c$ using \texttt{BreakMeshByBlockGenerator},
allowing for a displacement discontinuity. The jump is defined as
\begin{equation}
\jump{\bu} = \bu^{+} - \bu^{-} \quad \text{on } \Gamma_c.
\end{equation}

The unit normal vector is $\bn = (0,1)$, and the tangential direction is $\btan = (1,0)$.
The normal and tangential components of the jump are
\begin{equation}
\delta_n = \jump{\bu} \cdot \bn = \jump{u_y},
\qquad
\delta_t = \jump{\bu} \cdot \btan = \jump{u_x}.
\end{equation}

%=====================================================
\subsection{Bulk equations}

In each subdomain $\Omega_i$, $i=1,2$, we solve the static equilibrium equations
\begin{equation}
\nabla \cdot \bsigma = \bm{0} \quad \text{in } \Omega_i,
\end{equation}
with isotropic linear elasticity
\begin{equation}
\bsigma = \mathbb{C} : \bm{\varepsilon},
\qquad
\bm{\varepsilon} = \frac{1}{2}(\nabla \bu + \nabla \bu^{\top}).
\end{equation}

The material parameters are chosen as
\begin{equation}
E = 10^{9}, \qquad \nu = 0.3,
\end{equation}
such that the bulk stiffness is significantly larger than the interface stiffness.

%=====================================================
\subsection{Interface kinematics and equilibrium}

On $\Gamma_c$, the cohesive traction satisfies
\begin{equation}
\bsigma\,\bn = \btrac(\delta_n, \delta_t).
\end{equation}

Under mode-I loading, only the normal component is active,
\begin{equation}
t_n = \btrac \cdot \bn.
\end{equation}

%=====================================================
\subsection{Bilinear traction--separation law}

For monotonic mode-I loading, $\delta_t = 0$ and $\delta_n^{\max} = \delta_n$.
The bilinear traction law reduces to
\begin{equation}
t_n(\delta_n) =
\begin{cases}
K\,\delta_n, & 0 \le \delta_n \le \delta_n^0, \\[6pt]
(1-d)\,K\,\delta_n, & \delta_n^0 < \delta_n \le \delta_n^f, \\[6pt]
0, & \delta_n > \delta_n^f.
\end{cases}
\end{equation}

The damage variable is defined as
\begin{equation}
d =
\frac{\delta_n^f \left(\delta_n^{\max} - \delta_n^0\right)}
{\delta_n^{\max} \left(\delta_n^f - \delta_n^0\right)},
\qquad d \in [0,1].
\end{equation}

For monotonic loading, the softening branch can be written as
\begin{equation}
t_n(\delta_n) =
N \frac{\delta_n^f - \delta_n}{\delta_n^f - \delta_n^0},
\qquad
\delta_n^0 \le \delta_n \le \delta_n^f.
\end{equation}

The characteristic separations are
\begin{equation}
\delta_n^0 = \frac{N}{K},
\qquad
\delta_n^f = \frac{2 G_{Ic}}{N}.
\end{equation}

The fracture energy satisfies
\begin{equation}
G_{Ic} = \int_0^{\delta_n^f} t_n(\delta_n)\, d\delta_n
= \frac{1}{2} N \delta_n^f.
\end{equation}

%=====================================================
\subsection{Numerical parameters}

\begin{table}[h]
\centering
\small
\begin{tabular}{lll}
\toprule
Symbol & Meaning & Value \\
\midrule
$K$ & penalty stiffness & $5.0 \times 10^{3}$ \\
$N$ & normal strength & $50$ \\
$G_{Ic}$ & fracture energy & $0.5$ \\
$\delta_n^0$ & onset of softening & $1.0 \times 10^{-2}$ \\
$\delta_n^f$ & full decohesion & $2.0 \times 10^{-2}$ \\
\bottomrule
\end{tabular}
\end{table}

The penalty stiffness $K$ is chosen so that $\delta_n^0 = 0.01$ coincides
with the displacement at which the loading function $u_*(t)$ changes slope;
the elastic and softening branches then occupy comparable extents along
the $\delta_n$ axis.

%=====================================================
\subsection{Loading and boundary conditions}

The bottom boundary ($y=0$) is fixed, while the top boundary ($y=2$) is
displaced in the vertical direction according to
\begin{equation}
u_*(t) =
\begin{cases}
\dfrac{0.01}{0.1}\, t, & 0 \le t \le 0.1, \\[8pt]
0.01
+ \dfrac{0.03 - 0.01}{0.9}\,(t - 0.1),
& 0.1 < t \le 1.
\end{cases}
\end{equation}

The final displacement $u_*(1) = 0.03 = 1.5\,\delta_n^f$ ensures complete decohesion.

%=====================================================
\subsection{Weak form}

Find $\bu \in \mathcal{V}$ such that
\begin{equation}
\int_{\Omega_1 \cup \Omega_2}
\bsigma(\bu) : \nabla \bv \, dV
+
\int_{\Gamma_c}
\btrac(\jump{\bu}) \cdot \jump{\bv} \, dS
= 0
\end{equation}
for all admissible test functions $\bv$.

%=====================================================
\subsection{Verification criterion}

Let $(\delta_n^{(k)}, t_n^{(k)})$ denote simulation results. The solution is accepted if
\begin{equation}
\max_k \left| t_n^{(k)} - t_n^{\text{ref}}(\delta_n^{(k)}) \right| < 10^{-7},
\end{equation}
and
\begin{equation}
\max_k |t_t^{(k)}| < 10^{-10}.
\end{equation}

\end{document}
